\documentclass{article}
\usepackage[UTF8]{ctex}
\usepackage{amsmath,amssymb,amsfonts}
\usepackage{enumitem}
\usepackage{geometry}
\geometry{a4paper, margin=1in}
% === 错题本排版优化（不参与题目标识） ===
\setlength{\parindent}{0pt}
\setlength{\parskip}{0.35em plus 0.10em minus 0.08em}
\setlength{\emergencystretch}{3em}
\setlist{itemsep=0.25em, topsep=0.40em, parsep=0pt}
\allowdisplaybreaks[2]
% === 错题本排版优化结束 ===
\begin{document}
\section*{合肥工业大学2026年数学分析真题}
\section*{109．合肥工业大学2026年数学分析试卷}

\subsection{题目1}
一、（15 分）判断极限 \(\lim _{n \rightarrow \infty} \sum_{k=1}^{n} \ln \left(1+\frac{k}{n^{2}}\right)\) 是否存在？若存在，求出极限值．

\subsection{题目2}
二．（15 分）设 \(f(x)\) 在 \([0,1]\) 上连续，在 \((0,1)\) 上可导，\(f(0)=0, f(1)=1\) ，证明：在 \((0,1)\) 内存在不同的 \(\xi, \eta\) ，使得 \(f^{\prime}(\xi) f^{\prime}(\eta)=1\) ．

\subsection{题目3}
三．（15 分）设 \(f(x)\) 在 \([1,+\infty)\) 上连续，且 \(\lim _{x \rightarrow+\infty}[f(x)-3 \sqrt{x}]=7\) ．判断函数 \(f(x)\) 在 \([1,+\infty)\) 上是否一致连续，并说明你的理由。

\subsection{题目4}
四．（15 分）设 \(f(x)\) 在 \([a, b]\) 上有定义，并且在 \([a, b]\) 上每一点的极限都存在且为 0 ，证明：\(f(x)\) 在\([a, b]\) 上可积，并求 \(\int_{a}^{b} f(x) \mathrm{d} x\) ．

\subsection{题目5}
五．（15 分）设 \(f(x)\) 在 \([a, b]\) 上连续，且 \(\min _{x \in[a, b]} f(x)=3\) ，计算

\[
\lim _{n \rightarrow \infty}\left[\int_{a}^{b} \frac{1}{(f(x))^{n}} \mathrm{~d} x\right]^{\frac{1}{n}}
\]

\subsection{题目6}
六．（15 分）设 \(a_{n}=\int_{0}^{\frac{\pi}{4}} \tan ^{n} x \mathrm{~d} x\) ，计算 \(\sum_{n=1}^{\infty} \frac{a_{n}+a_{n+2}}{n}\) 的值，判断级数 \(\sum_{n=1}^{\infty} \frac{a_{n}}{\sqrt{n}}\) 是否收敛？并说明理由．

\subsection{题目7}
七．（15 分）设 \(F(x, y)\) 在 \(\left(x_{0}, y_{0}\right)\) 某邻域内有二阶连续的偏导数，且

\[
F\left(x_{0}, y_{0}\right)=0, F_{x}^{\prime}\left(x_{0}, y_{0}\right)=0, F_{x x}^{\prime \prime}\left(x_{0}, y_{0}\right) F_{y}^{\prime}\left(x_{0}, y_{0}\right) \neq 0
\]

证明：由方程 \(F(x, y)=0\) 在 \(\left(x_{0}, y_{0}\right)\) 某邻域内确定的隐函数 \(y=y(x)\) 在点 \(x_{0}\) 处取极值．

\subsection{题目8}
八．（15 分）求平面 \(x+2 y=1\) 与曲面 \(x^{2}+2 y^{2}+z^{2}=1\) 的交线上距离原点最近的点的坐标．

\subsection{题目9}
九．（15 分）求曲面 \(\left(x^{2}+y^{2}+z^{2}\right)^{3}=a^{3} x y z(a>0)\) 所围立体的体积．

附加题．（15 分）证明：\(F(\lambda)=\int_{0}^{+\infty} \frac{1-e^{-\lambda t}}{t} \cos t \mathrm{~d} t\) 在 \([0,+\infty)\) 上连续，在 \((0,+\infty)\) 内可导．
\end{document}
